\chapter{Kaposi Sarcoma}

\section*{Definition and Overview}
Kaposi sarcoma is a cancer that arises from the cells lining blood and lymph vessels, caused by infection with the human herpesvirus 8 (HHV-8), also called Kaposi sarcoma-associated herpesvirus. It appears as purple, red, or brown patches, plaques, or nodules on the skin, and can also affect the mouth, lymph nodes, and internal organs such as the lungs and gut. The condition occurs almost exclusively in people with weakened immunity, including people with HIV, organ transplant recipients, and older people of certain ethnic backgrounds. With antiretroviral therapy for HIV and treatment of the lesions, most people respond well, and prevention centres on controlling HIV and immunity.

\section*{Epidemiology}
Kaposi sarcoma is rare in Australia and most developed countries. It was common in people with untreated HIV before effective antiretroviral therapy, and its incidence has fallen dramatically since. It remains one of the most common cancers in people with HIV in sub-Saharan Africa and other regions where HIV is widespread. In Australia, most cases occur in people with HIV not on treatment, organ transplant recipients taking immunosuppressants, and older men of Mediterranean or Jewish background, in whom a chronic, slowly progressive form occurs.

\section*{Aetiology and Pathophysiology}
Kaposi sarcoma is caused by human herpesvirus 8, a virus transmitted through saliva, sexual contact, and blood. Most people infected with HHV-8 never develop the cancer; it develops when the immune system is weakened, allowing the virus to reactivate and drive the proliferation of infected endothelial cells into tumour tissue. The growths are abnormal blood vessel cells, which is why they appear purple and bleed easily. In people with HIV, immune recovery with antiretroviral therapy controls the virus and causes lesions to regress, while in transplant recipients, reducing immunosuppression can achieve the same. In the classic form in older people, the disease is slow-growing and often limited to the skin.

\section*{Clinical Features}
\begin{itemize}
    \item Purple, red, or brown patches, plaques, or nodules on the skin
    \item Lesions on the legs, feet, face, and oral cavity
    \item Lesions may be painless or painful and can swell or bleed
    \item Swelling of a limb from blockage of lymph channels
    \item Mouth lesions that bleed or interfere with eating
    \item Lung involvement: cough, breathlessness, and coughing blood
    \item Gut involvement: weight loss, abdominal pain, and bleeding
    \item Fever, night sweats, and weight loss with advanced disease
\end{itemize}

\section*{Differential Diagnosis}
The appearance of Kaposi sarcoma can be confused with other skin and systemic conditions.
\begin{itemize}
    \item \textbf{Bruises and blood collections:} dark lesions that change and fade
    \item \textbf{Bacterial or fungal skin infections:} in immunocompromised people
    \item \textbf{Other skin cancers:} melanoma, squamous cell, and basal cell carcinoma
    \item \textbf{Other vascular lesions:} haemangiomas and angiosarcoma
    \item \textbf{Other HIV-related conditions:} infections and lymphomas
    \item \textbf{Pigmented lesions:} moles and seborrhoeic keratoses
\end{itemize}

\section*{When to Seek Medical Care}
New purple or brown skin lesions, especially in people with HIV, transplant recipients, or those with immune suppression, should be assessed. Lesions that bleed, swell, or ulcerate, or symptoms such as cough, breathlessness, abdominal pain, or unexplained weight loss, warrant medical review. People with HIV should remain on antiretroviral therapy and attend regular reviews.

\textbf{Call 000 (or local emergency services) for:}
\begin{itemize}
    \item Coughing blood or severe breathlessness
    \item Heavy bleeding from a lesion
    \item Severe abdominal pain, vomiting blood, or black stools
    \item Collapse, confusion, or signs of severe infection
\end{itemize}

\section*{Diagnosis and Investigations}
\begin{itemize}
    \item \textbf{Skin biopsy:} the definitive diagnosis, with histology of the lesion
    \item \textbf{HIV testing:} recommended in anyone diagnosed with Kaposi sarcoma
    \item \textbf{HHV-8 testing:} antibody testing in selected cases
    \item \textbf{Chest imaging:} X-ray or CT for lung involvement
    \item \textbf{Endoscopy:} for gut involvement with bleeding or symptoms
    \item \textbf{CT scans:} for lymph node and internal disease
    \item \textbf{Assessment of immune function:} CD4 count in people with HIV
\end{itemize}

\section*{Management}
\begin{itemize}
    \item \textbf{Antiretroviral therapy:} the cornerstone for people with HIV, restoring immunity and controlling the virus
    \item \textbf{Reducing immunosuppression:} in transplant recipients, under specialist care
    \item \textbf{Local treatment:} topical gels, freezing, laser, or radiotherapy for skin lesions
    \item \textbf{Systemic chemotherapy:} for extensive or internal disease
    \item \textbf{Immunotherapy and targeted therapy:} newer options in selected cases
    \item \textbf{Symptom management:} treatment of swelling, bleeding, and pain
    \item \textbf{Monitoring:} regular review of lesions and response
    \item \textbf{Multidisciplinary care:} oncology, infectious disease, and dermatology
\end{itemize}

\section*{Complications}
\begin{itemize}
    \item Bleeding and infection of skin lesions
    \item Lymphoedema from blockage of lymphatic drainage
    \item Respiratory failure from lung involvement
    \item Bowel obstruction and bleeding from gut disease
    \item Effects of treatment
    \item Progression if immunity remains poor
\end{itemize}

\section*{Prognosis and Outcomes}
The outlook for Kaposi sarcoma depends on the state of the immune system and the extent of disease. In people with HIV, antiretroviral therapy causes most lesions to regress and dramatically improves survival, and Kaposi sarcoma is now uncommon in treated patients. Localised skin disease responds well to local treatment, and systemic therapy controls extensive disease. The key to prevention is maintaining immunity, whether through antiretroviral therapy or the careful management of immunosuppression.

\section*{Prevention and Self-Care}
\begin{itemize}
    \item If you have HIV, take antiretroviral therapy consistently and attend reviews
    \item Be aware of the signs of Kaposi sarcoma if you have HIV or take immunosuppressants
    \item Seek early assessment of new skin lesions
    \item Protect the skin and report bleeding, swelling, or pain
    \item Attend cancer and infectious disease follow-up
    \item Avoid further immune suppression where possible
    \item Access support services for HIV and cancer care
    \item Maintain overall health with nutrition, activity, and sleep
\end{itemize}

\section*{Sources and Further Reading}
The following reputable sources provide further information for healthcare students and patients seeking reliable, current advice about Kaposi sarcoma.
\begin{itemize}
    \item Cancer Council Australia. \textit{Kaposi sarcoma information.} https://www.cancer.org.au/
    \item Healthdirect Australia. \textit{Kaposi sarcoma.} https://www.healthdirect.gov.au/
    \item Positive Life NSW. \textit{HIV and health information.} https://positivelife.org.au/
    \item Australasian Society for HIV, Viral Hepatitis and Sexual Health Medicine. \textit{HIV care resources.} https://ashm.org.au/
    \item NHS. \textit{Kaposi's sarcoma.} https://www.nhs.uk/conditions/kaposis-sarcoma/
    \item American Cancer Society. \textit{Kaposi sarcoma.} https://www.cancer.org/
\end{itemize}
